\documentclass[12pt, a4paper]{article}

% -----------------------------------------------
% Packages
% -----------------------------------------------
\usepackage[T1]{fontenc}
\usepackage[utf8]{inputenc}
\usepackage{lmodern}          % Standard research-paper font (Latin Modern)
\usepackage[margin=1in]{geometry}
\usepackage{amsmath, amssymb}
\usepackage{graphicx}
\usepackage{booktabs}
\usepackage{array}
\usepackage{tabularx}
\usepackage{longtable}
\usepackage{float}
\usepackage{hyperref}
\usepackage{url}
\usepackage{caption}
\usepackage{subcaption}
\usepackage{setspace}
\usepackage{parskip}
\usepackage{titlesec}
\usepackage{fancyhdr}
\usepackage{xcolor}
\usepackage{microtype}
\usepackage{enumitem}
\usepackage{natbib}

% -----------------------------------------------
% Page style
% -----------------------------------------------
\pagestyle{fancy}
\fancyhf{}
\rhead{Data Analysis 1 --- Spring 2026}
\lhead{Coffee Quality Tier Prediction}
\cfoot{\thepage}

% -----------------------------------------------
% Section formatting
% -----------------------------------------------
\titleformat{\section}{\large\bfseries}{}{0em}{\Roman{section}.\quad}
\titleformat{\subsection}{\normalsize\bfseries}{}{0em}{\arabic{section}.\arabic{subsection}\quad}

% -----------------------------------------------
% Misc helpers
% -----------------------------------------------
\newcommand{\placeholder}[1]{%
  \begin{center}
    \fbox{\parbox{0.85\textwidth}{\centering\vspace{1.5cm}
      \textit{[Placeholder: #1]}\vspace{1.5cm}}}
  \end{center}
}

% -----------------------------------------------
% Title block
% -----------------------------------------------
\begin{document}

\begin{titlepage}
  \centering
  \vspace*{1cm}

  % University logo placeholder
  \placeholder{University Logo --- American University of Phnom Penh}

  \vspace{1.5cm}
  {\LARGE\bfseries A Logistic Regression Model: Predicting\\[0.4em]
   Coffee Quality Tier without Cupping\par}

  \vspace{1.5cm}
  {\large\textbf{Course:} Data Analysis 1\par}
  \vspace{0.4cm}
  {\large\textbf{Lecturer:} Dr.\ HAS Sothea\par}

  \vspace{1.2cm}
  {\large\textbf{Semester:} Spring 2026\par}

  \vspace{1cm}
  {\large\textbf{Team Members:}\par}
  \vspace{0.4cm}
  {\large Khan Mengheap --- 2024497\par}
  \vspace{0.2cm}
  {\large Taing Muyleang --- 2024355\par}

  \vfill
  {\large\textbf{Date:} 26 April 2026\par}
\end{titlepage}

% -----------------------------------------------
\setcounter{section}{0}

% -----------------------------------------------
\section{Introduction}
% -----------------------------------------------

Coffee is one of the most widely consumed beverages and agricultural exports for many developing
countries. A coffee's quality grade (Commercial, Premium, or Specialty) has significant economic
consequences for farmers, exporters, importers, and roasters. Being able to predict coffee quality
early in the supply chain can help stakeholders make better sourcing decisions, reduce waste, and
reward high-quality producers with appropriate pricing.

This project aims to develop and evaluate multiclass classification models that predict the quality
tier of Arabica coffee samples using two types of features:
\begin{enumerate}
  \item Production-level attributes such as growing altitude, processing method, bean variety, and
        physical defects.
  \item Sensory evaluation scores assigned by certified Q-graders during cupping sessions.
\end{enumerate}

By comparing these two feature sets, we aim to determine how well coffee quality tier can be
predicted from production-side features alone, without any cupping data, and to explore the
implications for real-world applications such as pre-screening, pricing, and farmer self-assessment.

The dataset used in the project is the Coffee Quality Institute (CQI) Arabica Quality Database
\citep{LeDoux2018}, a public dataset that contains 1,311 professionally evaluated Arabica coffee
samples from 34 countries, spanning harvest years from 2009 to 2018. Each record includes origin
information, processing details, physical defect counts, altitude data, and 10 individual sensory
scores (Aroma, Flavor, Aftertaste, Acidity, Body, Balance, Uniformity, Clean Cup, Sweetness, and
Cupper Points), each scored on a 1--10 scale by certified Q-graders, which together sum to the Total
Cup Points score (maximum 100). From Total Cup Points, we derived our target variable
\texttt{quality\_tier} using the Specialty Coffee Association's standard grading threshold:
Commercial (below 80), Premium (80--84), and Specialty (85 and above) \citep{SCA2005}.

After preprocessing---which primarily includes parsing messy altitude and harvest year fields,
handling missing values, and dropping a small number of rows with unparseable data---this dataset
contains 1,251 Arabica coffee samples across three different quality tiers: Commercial (13.9\%),
Premium (78.0\%), and Specialty (8.1\%).

% -----------------------------------------------
\section{Data Preprocessing}
% -----------------------------------------------

The raw CQI dataset contains 1,311 rows and 44 columns mixing sensory scores, production
metadata, administrative identifiers, and pre-computed cleanup fields of uneven quality. This
section documents every step of our cleaning process in the order it was performed in the notebook.

\subsection{Initial Data Inspection}

The dataset was loaded from the CQI GitHub repository using \texttt{pd.read\_csv()}, returning
1,311 rows and 44 columns. All column names were converted to \texttt{snake\_case} format to
ensure consistency throughout the analysis. A column-level summary was then computed using
\texttt{df.dtypes}, \texttt{df.nunique()}, and \texttt{df.isnull().sum()} to capture data
types, cardinality, and missingness for every column.

Calling \texttt{df.duplicated().sum()} confirmed that there were no duplicate rows. However,
many columns had missing values, which exhibited a clear pattern: the columns with the highest
proportion of missing values were administrative identifiers, namely \texttt{lot\_number}
(79.4\% missing), \texttt{farm\_name} (27.2\%), \texttt{mill} (23.7\%), \texttt{producer}
(17.5\%), and \texttt{company} (15.9\%). This pattern is consistent with the CQI submission
process, where producers may submit coffees without complete administrative paperwork.

Among coffee-attribute columns, missingness appeared in \texttt{color} (20.4\%), the three
altitude-related fields (each 17.3\% missing), the raw altitude column (17.0\%),
\texttt{variety} (15.3\%), and \texttt{processing\_method} (11.6\%). Each of these required
individual handling decisions in subsequent preprocessing steps.

Most data types were appropriate (the ten sensory scores stored as \texttt{float64}, defect
counts as \texttt{int64}), but two columns stored as \texttt{object} required custom parsing:
\texttt{altitude} (mixed unit formats, e.g., ``1500 m'', ``4200 feet'', ``1000--1500m''), and
\texttt{harvest\_year} (inconsistent formatting with year ranges, quarter notation, and
multilingual month names). The pre-computed altitude cleanup columns
(\texttt{altitude\_low\_meters}, \texttt{altitude\_high\_meters}, and
\texttt{altitude\_mean\_meters}) appeared clean as \texttt{float64} but, as discovered later in
Section~II.3, contained systematic parsing bugs.

\subsection{Drop Irrelevant Columns}

Based on the inspection in the previous section, 23 columns were identified as not useful for
predicting coffee quality and were removed before further analysis. These fell into four categories:

\textbf{Zero-variance column.} The \texttt{species} column contained only one unique value
``Arabica'' for all rows. Since the dataset is restricted to Arabica coffee, this column is useless
and was dropped.

\textbf{Administrative identifiers.} Thirteen columns recorded paperwork metadata:
\texttt{unnamed:\_0} (a stale row-index column), \texttt{lot\_number}, \texttt{ico\_number},
\texttt{farm\_name}, \texttt{mill}, \texttt{producer}, \texttt{company}, \texttt{owner},
\texttt{owner\_1}, \texttt{in\_country\_partner}, \texttt{certification\_body},
\texttt{certification\_address}, and \texttt{certification\_contact}. Including these columns
in the model would force it to memorize specific producers or certification bodies rather than
learning generalizable patterns about coffee quality.

\textbf{Process metadata.} Four columns recorded certification logistics or batch information rather
than coffee properties: \texttt{grading\_date} and \texttt{expiration} describe when each sample
was processed, while \texttt{bag\_weight} and \texttt{number\_of\_bags} describe the size of
the submission batch.

\textbf{Redundant or buggy columns.} Five columns either duplicated information or contained
errors. The four altitude-related columns (\texttt{unit\_of\_measurement},
\texttt{altitude\_low\_meters}, \texttt{altitude\_high\_meters}, and
\texttt{altitude\_mean\_meters}) were the attempted cleanup of the raw altitude text but contained
systematic parsing bugs (see Section~II.3). The \texttt{region} column had 343 unique country
region names with most having only 1--2 samples; geographic information was instead captured
through \texttt{country\_of\_origin}, later regrouped into broader continental regions in
Section~II.8.

After this drop, the working dataset contained 1,311 rows and 21 columns.

\subsection{Initial Processing Decisions}

After dropping irrelevant columns, missingness was re-examined across the remaining 21 columns.
The missing pattern persisted: \texttt{color} (267 missing, 20.4\%), raw altitude (223 missing,
17.0\%), \texttt{variety} (201 missing, 15.3\%), \texttt{processing\_method} (152 missing,
11.6\%), and \texttt{harvest\_year} (47 missing, 3.6\%), with single missing values in
\texttt{country\_of\_origin} and \texttt{quakers}.

One row with \texttt{total\_cup\_points == 0} was dropped as a data-entry error (the practical
minimum cupping score is around 60), leaving 1,310 rows.

The classification target \texttt{quality\_tier} was constructed from \texttt{total\_cup\_points}
using the SCA grading thresholds \citep{SCA2005}:
\begin{itemize}
  \item \textbf{Commercial:} total cup points below 80
  \item \textbf{Premium:} total cup points between 80 and 84.99
  \item \textbf{Specialty:} total cup points 85 or above
\end{itemize}

Applied to the data, the target distribution was 174 Commercial samples (13.3\%), 1,033 Premium
samples (78.9\%), and 103 Specialty samples (7.9\%). This class imbalance shaped several decisions
in model development (Section~IV).

\subsection{Altitude Parser}

The \texttt{altitude} column required the most preprocessing effort. The four altitude cleanup
columns provided in the dataset contained systematic bugs that made them unreliable, so a
customized parser was developed using the raw altitude text column.

Inspection of \texttt{altitude\_mean\_meters} revealed impossible values: the maximum was
190,164~m (for reference, Mount Everest's summit is 8,849~m). Two problem categories were
identified:

\begin{itemize}
  \item \textbf{Decimal-stripping errors:} Row 1144 had a raw altitude of ``1901.64'' (Guatemala)
        but was stored as 190,164.0; row 1040 had ``1100.00 mosl'' (Nicaragua) stored as 110,000.0.
        In both cases, the parser stripped the decimal point and treated the digits as a single integer.
  \item \textbf{Producer-side reporting errors:} Several Guatemalan rows reported ``3280'' with
        \texttt{unit\_of\_measurement = `m'}, likely confusing feet with meters; multiple Myanmar
        rows reported altitudes between 3,800 and 4,287~m, which are biologically impossible; one
        Brazilian row reported ``11000 metros''.
\end{itemize}

A \texttt{parse\_altitude()} function was written with four steps:
\begin{enumerate}
  \item Extract all numeric values using the regular expression \verb|\d+\.?\d*| (capturing both
        integers and decimals, preventing the decimal-stripping bug).
  \item If two numbers are extracted (e.g., ``1000--1500m''), use the midpoint.
  \item If the string contains \textit{feet}, \textit{ft}, or an apostrophe, convert from feet to
        meters using the factor 0.3048.
  \item If the resulting altitude falls outside the biologically plausible range of 100--3000~m, flag
        the value as missing.
\end{enumerate}

After parsing, 948 rows had valid altitudes and 360 rows were flagged. Missing altitudes were
imputed with the median altitude of valid samples from the same \texttt{country\_of\_origin}
to preserve real geographic signal. Two rows whose \texttt{country\_of\_origin} was missing
could not be imputed and were dropped, leaving 1,308 rows.

The final \texttt{altitude\_imputed\_m} column has a mean of 1,301~m, median of 1,350~m, and
standard deviation of 398~m, consistent with global Arabica production.

\subsection{Harvest Year Parser}

The \texttt{harvest\_year} column was stored as \texttt{object} because of inconsistent
submission formats. The 46 unique values fell into five formatting patterns: plain years
(e.g., ``2012''), year ranges (e.g., ``2013/2014''), two-digit year codes (e.g., ``08/09 crop''),
quarter notation (e.g., ``1T/2011''), and date text in multiple languages (English and Spanish),
including junk entries such as ``TEST'' and ``mmm''.

A \texttt{parse\_harvest\_year()} function was written using two regex-based strategies:
\begin{enumerate}
  \item Find any 4-digit year in the range 2000--2029 using \verb|\b(20[0-2]\d)\b|. If multiple
        years are found, take the later year.
  \item Fallback for two-digit codes: search for 2-digit codes following a slash using
        \verb|/(\d{2})\b| (e.g., ``4T/10'' $\to$ 2010).
  \item Strings with no extractable year information return \texttt{NaN}.
\end{enumerate}

The parser successfully extracted years for 1,251 of 1,308 rows. The 57 unparseable rows were
dropped. The cleaned \texttt{harvest\_year\_clean} column spans 2009 to 2018 with a median of
2014 and standard deviation of 1.85 years.

\subsection{Variety Grouping}

The \texttt{variety} column contained 29 unique values plus 156 missing entries (12.5\%
missingness). To reduce cardinality while preserving meaningful information, varieties were grouped
into 10 botanical families using the Arabica variety taxonomy maintained by World Coffee Research
and the Specialty Coffee Association:

\begin{itemize}
  \item \textbf{Bourbon:} Bourbon and Yellow Bourbon.
  \item \textbf{Typica:} Typica, Mundo Novo, Blue Mountain, and Java.
  \item \textbf{SL\_Kenyan:} SL14, SL28, and SL34.
  \item \textbf{Specialty:} Geisha, Pacamara, Pacas, Marigojipe, Peaberry, and Moka Peaberry
        (rare specialty varieties, each with fewer than 15 samples).
  \item \textbf{Caturra, Catuai, Catimor, Hawaiian Kona:} Kept as their own groups (each with at
        least 20 samples and a distinct cultivation identity).
  \item \textbf{Other:} Remaining low-sample varieties.
  \item \textbf{Unknown:} The 156 missing entries, preserved as their own informative category.
\end{itemize}

The final distribution: Bourbon (259), Caturra (252), Typica (247), Unknown (156), Other (127),
Catuai (72), Hawaiian Kona (44), SL\_Kenyan (38), Specialty (36), and Catimor (20).

\subsection{Processing Method Cleanup}

The \texttt{processing\_method} column had five unique values plus 100 missing entries
(8.0\% missingness). Slash-notation labels were shortened (e.g., ``Washed / Wet'' $\to$
``Washed''), and the 12-sample ``Pulped natural / honey'' category was merged into
``Semi-washed'' due to process similarity. Missing entries were labelled ``Unknown''.

Final distribution: Washed (808, 64.6\%), Natural (249, 19.9\%), Unknown (100, 8.0\%),
Semi-washed (68, 5.4\%), and Other (26, 2.1\%).

\subsection{Color Cleanup}

The \texttt{color} column contained three unique values plus 220 missing entries (17.6\%
missingness): Green (839, 67.1\%), Bluish-Green (112, 9.0\%), and Blue-Green (80, 6.4\%).
The distinction between ``Bluish-Green'' and ``Blue-Green'' is linguistic rather than meaningful;
both describe freshly-milled, high-moisture green coffee beans. They were merged into a single
``Blue-Green'' category. Missing entries were labelled ``Unknown''.

Final distribution: Green (839, 67.1\%), Unknown (220, 17.6\%), Blue-Green (192, 15.4\%).

\subsection{Country to Region Grouping}

The \texttt{country\_of\_origin} column contained 34 unique values with severe imbalance
(the top five countries accounted for 62\% of samples, while 15 countries had fewer than 10
samples each). Countries were grouped into five coffee-growing regions based on the International
Coffee Organization's geographic classifications:

\begin{itemize}
  \item \textbf{Central America:} Mexico, Guatemala, Honduras, Costa Rica, Nicaragua,
        El Salvador, Panama.
  \item \textbf{South America:} Colombia, Brazil, Peru, Ecuador.
  \item \textbf{East Africa:} Ethiopia, Tanzania, Uganda, Kenya, Malawi, Burundi, Rwanda,
        Zambia, and C\^{o}te d'Ivoire.
  \item \textbf{Southeast Asia:} Taiwan, Thailand, Indonesia, China, Myanmar, Vietnam,
        Philippines, Laos, and Papua New Guinea.
  \item \textbf{Caribbean/US:} Hawaii, Haiti, Puerto Rico, U.S. mainland, and Mauritius.
\end{itemize}

Final distribution: Central America (561, 44.8\%), South America (309, 24.7\%), Southeast Asia
(161, 12.9\%), East Africa (147, 11.8\%), and Caribbean/US (73, 5.8\%).

\subsection{Moisture Cleanup}

The \texttt{moisture} column appeared complete but 229 rows (18.3\%) had a moisture reading of
exactly 0, which is physically impossible for coffee beans (target range: 10--12\% for shelf
stability). These were interpreted as hidden missingness.

Regional medians for moisture (computed from non-zero values) had a spread of only 0.01
percentage point across all five regions, so the simpler global median of 0.110 (11.0\%) was used
to fill all 229 zero values.

After imputation: minimum 0.01, median 0.11, maximum 0.28, standard deviation 0.018.

\subsection{Quakers Missing Value}

The \texttt{quakers} column contained a single missing value, imputed with the mode 0, since
over 92\% of samples were already at zero. Its effect on summary statistics and model development
is negligible at a sample size of 1,251.

\subsection{Drop \texttt{color\_group}}

Chi-square tests of independence between each cleaned categorical feature and
\texttt{quality\_tier} (full results in Section~III.8) yielded the following ranking:
\texttt{region\_group} (188.98) $>$ \texttt{variety\_group} (126.41) $>$
\texttt{processing\_group} (47.78) $>$ \texttt{color\_group} (42.08).

\texttt{color\_group} showed the weakest association, the highest missingness of any retained
categorical feature (17.6\%), and only three categories (one being a missing-data placeholder).
It was therefore dropped from the feature set. The working dataset after this step contained
\textbf{1,251 rows and 20 columns}.

% -----------------------------------------------
\section{Exploratory Data Analysis}
% -----------------------------------------------

This section explores the cleaned dataset to understand how production features and sensory scores
relate to the quality tier.

\subsection{Target Distribution}

The dataset is highly imbalanced:
\begin{itemize}
  \item Premium: 78.0\%
  \item Commercial: 13.9\%
  \item Specialty: 8.1\%
\end{itemize}

To address this imbalance, \texttt{class\_weight=`balanced'} was applied in one of the logistic
regression models (Model~A -- balanced) to ensure better performance across all classes.

\subsection{Numeric Production Features}

\begin{table}[H]
  \centering
  \caption{Descriptive Statistics for Numeric Production Features}
  \begin{tabular}{lrrrrrrr}
    \toprule
    \textbf{Feature} & \textbf{Count} & \textbf{Mean} & \textbf{SD}
      & \textbf{Min} & \textbf{25\%} & \textbf{Median} & \textbf{Max} \\
    \midrule
    altitude\_imputed\_m   & 1,251 & 1,310.7 & 390.0 & 100.0  & 1,054.5 & 1,350.0 & 2,560.0 \\
    harvest\_year\_clean   & 1,251 & 2,013.7 &   1.8 & 2,009  & 2,012   & 2,014   & 2,018   \\
    moisture               & 1,251 &   0.111 & 0.018 &   0.01 &   0.11  &   0.11  &   0.28  \\
    category\_one\_defects & 1,251 &    0.42 &  1.84 &   0    &   0     &   0     &  31     \\
    category\_two\_defects & 1,251 &    3.66 &  5.41 &   0    &   0     &   2     &  55     \\
    quakers                & 1,251 &    0.17 &  0.80 &   0    &   0     &   0     &  11     \\
    \bottomrule
  \end{tabular}
\end{table}

\begin{figure}[H]
  \centering
  \placeholder{Figure 1: Distribution of Numeric Production Features}
  \caption{Distribution of Numeric Production Features}
\end{figure}

Key observations:
\begin{itemize}
  \item \textbf{Altitude} (\texttt{altitude\_imputed\_m}): The average altitude is 1,311~m, with
        most samples between 1,054~m and 1,600~m.
  \item \textbf{Harvest year} (\texttt{harvest\_year\_clean}): Samples range from 2009 to 2018,
        with most collected between 2012 and 2015.
  \item \textbf{Moisture:} The average moisture level is about 11.1\%, consistent with standard
        industry levels.
  \item \textbf{Category one defects:} Most coffees have no primary defects (median = 0); the
        distribution is heavily right-skewed (maximum = 31).
  \item \textbf{Category two defects:} Similarly right-skewed (maximum = 55, average 3.66).
  \item \textbf{Quakers:} Most samples have zero quakers (median = 0, mean = 0.17).
        Due to low variation and weak correlation with the target, \texttt{quakers} and
        \texttt{harvest\_year\_clean} were excluded from the final model.
\end{itemize}

\subsection{Sensory Scores}

\begin{table}[H]
  \centering
  \caption{Descriptive Statistics for Sensory Score Features}
  \begin{tabular}{lrrrrrrr}
    \toprule
    \textbf{Feature} & \textbf{Count} & \textbf{Mean} & \textbf{SD}
      & \textbf{Min} & \textbf{25\%} & \textbf{Median} & \textbf{Max} \\
    \midrule
    aroma         & 1,251 & 7.57 & 0.32 & 5.08 &  7.42 &  7.58 &  8.75 \\
    flavor        & 1,251 & 7.52 & 0.34 & 6.08 &  7.33 &  7.50 &  8.83 \\
    aftertaste    & 1,251 & 7.40 & 0.35 & 6.17 &  7.25 &  7.42 &  8.67 \\
    acidity       & 1,251 & 7.54 & 0.32 & 5.25 &  7.33 &  7.50 &  8.75 \\
    body          & 1,251 & 7.52 & 0.29 & 5.25 &  7.33 &  7.50 &  8.58 \\
    balance       & 1,251 & 7.52 & 0.35 & 6.08 &  7.33 &  7.50 &  8.75 \\
    uniformity    & 1,251 & 9.85 & 0.49 & 6.00 & 10.00 & 10.00 & 10.00 \\
    clean\_cup    & 1,251 & 9.84 & 0.74 & 0.00 & 10.00 & 10.00 & 10.00 \\
    sweetness     & 1,251 & 9.91 & 0.45 & 1.33 & 10.00 & 10.00 & 10.00 \\
    cupper\_points& 1,251 & 7.50 & 0.43 & 5.17 &  7.25 &  7.50 &  8.67 \\
    \bottomrule
  \end{tabular}
\end{table}

\begin{figure}[H]
  \centering
  \placeholder{Figure 2: Distribution of Numeric Sensory Features}
  \caption{Distribution of Numeric Sensory Features}
\end{figure}

The ten sensory scores divide into two groups:
\begin{itemize}
  \item \textbf{Analytical scores} (aroma, flavor, aftertaste, acidity, body, balance,
        cupper\_points): mostly fall between 7.0 and 8.5, with standard deviations from 0.29 to
        0.43. These scores vary across samples, making them useful for distinguishing quality levels.
  \item \textbf{Protocol scores} (uniformity, clean\_cup, sweetness): generally very high, often
        close to 10, with mean values above 9.85 and standard deviations between 0.45 and 0.74.
        Since most samples receive similarly high scores, these features are less useful for
        distinguishing quality levels.
\end{itemize}

\subsection{Categorical Features}

\begin{figure}[H]
  \centering
  \placeholder{Figure 3: Distribution of Categorical Features}
  \caption{Distribution of Categorical Features}
\end{figure}

\begin{itemize}
  \item \textbf{Variety} (\texttt{variety\_group}): Bourbon (259), Caturra (252), and Typica (247)
        make up around 60\% of the dataset.
  \item \textbf{Processing method} (\texttt{processing\_group}): Washed (808) dominates; Natural
        (249), Unknown (100), Semi-washed (68), and Other (26).
  \item \textbf{Region} (\texttt{region\_group}): Central America (561), South America (309),
        Southeast Asia (161), East Africa (147), and Caribbean/US (73).
\end{itemize}

\subsection{Numeric Features vs.\ Quality Tier}

\begin{figure}[H]
  \centering
  \placeholder{Figure 4: Numeric Features vs. Quality Tier}
  \caption{Numeric Production Features vs.\ Quality Tier (Boxplots)}
\end{figure}

The boxplots show that:
\begin{itemize}
  \item \textbf{Altitude:} Altitude increases as coffee quality improves; Specialty is grown at
        higher altitudes, consistent with the view that altitude slows bean growth and improves
        flavor development \citep{Wintgens2009}.
  \item \textbf{Category one defects:} Commercial coffee shows a wider distribution; Premium and
        Specialty coffees have very low and similar defect levels.
  \item \textbf{Category two defects:} Commercial coffee shows higher defect levels; the difference
        between groups is clearer than for category one defects.
  \item \textbf{Moisture, harvest year, and quakers:} All quality groups show very similar
        distributions with no clear differences, confirming these features do not help in predicting
        quality.
\end{itemize}

\subsection{Categorical Features vs.\ Quality Tier}

\begin{figure}[H]
  \centering
  \placeholder{Figure 5: Categorical Features vs. Quality Tier}
  \caption{Categorical Features vs.\ Quality Tier (Grouped Bar Plots)}
\end{figure}

The grouped bar plots show:
\begin{itemize}
  \item \textbf{Variety:} SL\_Kenyan has the highest proportion of Specialty coffee; Typica has
        the lowest Specialty rate and the highest Commercial share.
  \item \textbf{Processing method:} Natural-processed coffees tend to produce more Specialty
        coffee, possibly because natural processing enhances flavor complexity. Semi-washed coffees
        are mostly Premium.
  \item \textbf{Region:} East Africa has the highest proportion of Specialty coffee. Caribbean/US
        also shows a relatively high Specialty rate. Southeast Asia has the lowest Specialty rate
        and a higher share of Commercial coffee.
\end{itemize}

\subsection{Correlation Heatmap}

\begin{figure}[H]
  \centering
  \placeholder{Figure 6: Correlation Heatmap -- Pearson}
  \caption{Correlation Heatmap -- Pearson}
\end{figure}

\begin{figure}[H]
  \centering
  \placeholder{Figure 7: Correlation Heatmap -- Spearman}
  \caption{Correlation Heatmap -- Spearman}
\end{figure}

Key findings from the Pearson correlation heatmap:
\begin{itemize}
  \item \textbf{Sensory scores are strongly correlated:} The strongest pairwise correlations are
        flavor--aftertaste ($r = 0.86$), aftertaste--balance ($r = 0.77$), and aroma--flavor
        ($r = 0.75$). To avoid multicollinearity, Model~B uses only \texttt{flavor} as a
        representative feature.
  \item \textbf{Flavor is the strongest predictor of \texttt{total\_cup\_points}:} Flavor
        ($r = 0.83$), aftertaste ($r = 0.82$), and balance ($r = 0.78$) have the highest
        correlations with the target. Protocol scores (uniformity, clean\_cup, sweetness) show
        weaker relationships ($r = 0.44$--$0.61$) and were excluded from the model.
  \item \textbf{Production features show weak linear relationships:} The strongest
        production-related correlations are \texttt{category\_two\_defects} ($r = -0.29$) and
        \texttt{altitude} ($r = 0.21$). Other production features have little to no linear
        correlation with the target.
\end{itemize}

Spearman results are broadly consistent with Pearson. Notable differences: altitude rises from 0.21
to 0.27 (suggesting mild non-linearity), while \texttt{category\_two\_defects} drops in magnitude
from $-0.29$ to $-0.17$ (suggesting Pearson was inflated by extreme defect outliers). Harvest year
and quakers show near-zero correlations under both methods, confirming they should be excluded.

\subsection{Chi-Squared Tests}

Chi-square tests were used to assess the association between each cleaned categorical feature and
\texttt{quality\_tier}.

\begin{table}[H]
  \centering
  \caption{Chi-Square Test Results for Categorical Features vs.\ Quality Tier}
  \begin{tabular}{lrrl}
    \toprule
    \textbf{Feature} & $\boldsymbol{\chi^2}$ & \textbf{df} & \textbf{p-value} \\
    \midrule
    region\_group     & 188.98 & 8 & $1.34 \times 10^{-36}$ *** \\
    variety\_group    & 126.41 & 18 & $2.57 \times 10^{-18}$ *** \\
    processing\_group &  47.78 &  8 & $1.09 \times 10^{-7}$  *** \\
    color\_group      &  42.08 &  4 & $1.60 \times 10^{-8}$  *** \\
    \bottomrule
    \multicolumn{4}{l}{\footnotesize *** $p < 0.001$}
  \end{tabular}
\end{table}

All four features showed statistically significant associations at $p < 0.001$. The chi-square
statistic was used as a relative ranking of association strength:
\begin{itemize}
  \item \texttt{region\_group} is the strongest predictor, consistent with the visual pattern that
        East Africa shows the highest Specialty rate and Southeast Asia the lowest.
  \item \texttt{variety\_group} also shows a strong association, matching the evidence that
        SL\_Kenyan varieties dominate the Specialty tier.
  \item \texttt{processing\_group} shows a weaker but significant association.
  \item \texttt{color\_group} has the weakest association and was removed from the feature set
        (see Section~II.11).
\end{itemize}

% -----------------------------------------------
\section{Model Development}
% -----------------------------------------------

Based on the EDA findings, logistic regression was selected as the modelling algorithm. It is
interpretable, appropriate for multi-class classification, and well-suited to the relatively small,
mixed-type feature set. Four models were built to allow structured comparisons.

\subsection{Data Splitting}

The dataset of 1,251 rows was split into 80\% training (1,000 rows) and 20\% test (251 rows) sets
using stratified sampling (\texttt{stratify = y}) to maintain class proportions, with
\texttt{random\_state = 42} for reproducibility.

\begin{table}[H]
  \centering
  \caption{Train/Test Split with Class Distribution}
  \begin{tabular}{lrrrr}
    \toprule
    \textbf{Split} & \textbf{Total} & \textbf{Commercial} & \textbf{Premium} & \textbf{Specialty} \\
    \midrule
    Training (80\%) & 1,000 & 139 (13.9\%) & 780 (78.0\%) & 81 (8.1\%) \\
    Test (20\%)     &   251 &  35 (13.9\%) & 196 (78.1\%) & 20 (8.0\%) \\
    \bottomrule
  \end{tabular}
\end{table}

\subsection{Model Selection \& Training}

Four models were trained:

\begin{itemize}
  \item \textbf{Dummy Classifier:} Used as a baseline by always predicting ``Premium''. Achieved
        78.1\% accuracy but a macro-F1 of 0.29, confirming it completely fails on Commercial and
        Specialty classes.
  \item \textbf{Model A -- Unbalanced:} Logistic regression on production features with default
        settings, tending to favor the dominant Premium class.
  \item \textbf{Model A -- Balanced:} Logistic regression on production features with
        \texttt{class\_weight=`balanced'}, giving more weight to minority classes. This was
        selected as the primary Model~A due to better minority-class performance.
  \item \textbf{Model B -- Flavor Only:} Logistic regression trained on the single \texttt{flavor}
        score (Pearson $r = 0.83$ with \texttt{total\_cup\_points}). Serves as a benchmark for
        sensory-based prediction.
\end{itemize}

% -----------------------------------------------
\section{Results and Evaluation}
% -----------------------------------------------

\subsection{Performance Metrics}

All models were evaluated on Accuracy, Macro-F1, Precision, and Recall. Due to class imbalance,
Macro-F1 is the primary metric, as it assigns equal weight to all classes and prevents the majority
class (Premium) from dominating the evaluation.

\begin{table}[H]
  \centering
  \caption{Model Performance Comparison}
  \begin{tabular}{lrrrr}
    \toprule
    \textbf{Model} & \textbf{Accuracy} & \textbf{Macro F1} & \textbf{Precision} & \textbf{Recall} \\
    \midrule
    Dummy Classifier   & 0.781 & 0.292 & 0.260 & 0.333 \\
    Model A -- Unbalanced & 0.785 & 0.409 & 0.615 & 0.397 \\
    Model A -- Balanced   & 0.526 & 0.463 & 0.475 & 0.644 \\
    Model B -- Flavor Only& 0.912 & 0.832 & 0.909 & 0.788 \\
    \bottomrule
  \end{tabular}
\end{table}

\subsection{Model Comparison}

\begin{itemize}
  \item \textbf{Dummy Classifier:} High accuracy (78.1\%) but very low Macro-F1 (0.292) confirms
        that accuracy alone is misleading under class imbalance.
  \item \textbf{Model A -- Unbalanced:} Slightly better than the dummy baseline (Macro-F1 = 0.409),
        but still favors the dominant Premium class. This suggests that production features alone
        provide limited predictive power when class imbalance is not addressed.
  \item \textbf{Model A -- Balanced:} Accuracy decreases to 52.6\% while Macro-F1 improves to
        0.463 and Recall rises to 0.644, indicating better detection of Commercial and Specialty
        samples at the expense of overall accuracy.
  \item \textbf{Model B -- Flavor Only:} Best overall performance: 91.2\% accuracy, Macro-F1 of
        0.832, Precision 0.909, and Recall 0.788. A single sensory feature substantially
        outperforms all production-feature models.
\end{itemize}

\subsection{Confusion Matrix Analysis}

\begin{figure}[H]
  \centering
  \placeholder{Figure 8: Confusion Matrices for All Four Models}
  \caption{Confusion Matrices}
\end{figure}

\begin{itemize}
  \item \textbf{Dummy:} Predicts every sample as ``Premium''; completely fails on Commercial and
        Specialty.
  \item \textbf{Model A -- Unbalanced:} Correctly identifies most Premium samples but misclassifies
        nearly all Specialty and Commercial samples.
  \item \textbf{Model A -- Balanced:} Identifies more Specialty coffee but still struggles with
        Commercial due to low precision.
  \item \textbf{Model B:} Mostly correct predictions with only slight confusion between Premium and
        Specialty.
\end{itemize}

\subsection{Feature Importance}

\begin{figure}[H]
  \centering
  \placeholder{Figure 9: Feature Importance for Model A Balanced (Specialty Class)}
  \caption{Feature Importance for Model~A Balanced}
\end{figure}

Coefficients for the Specialty class from the balanced logistic regression model:

\textbf{Strongest positive predictors:}
\begin{itemize}
  \item \texttt{processing\_group\_Semi-washed} ($+0.74$): strongest positive predictor of
        Specialty classification.
  \item \texttt{altitude\_imputed\_m} ($+0.57$): higher altitude increases the probability of
        Specialty classification.
  \item \texttt{variety\_group\_Caturra} ($+0.46$) and \texttt{variety\_group\_Catimor}
        ($+0.35$): both positively associated with Specialty relative to the Bourbon reference.
\end{itemize}

Notably, \texttt{region\_group\_East\_Africa} and \texttt{variety\_group\_SL\_Kenyan} do not
show strong positive coefficients despite their visual dominance in the bivariate plots. This is due
to multicollinearity: most SL\_Kenyan samples originate from East African countries, so the model
splits the predictive signal between the two related dummy variables. Additionally, the Caribbean/US
reference region has a high Specialty rate (driven by Hawaiian Kona), which reduces the apparent
East Africa coefficient.

\textbf{Strongest negative predictors:}
\texttt{region\_group\_Central\_America} ($-1.70$) and
\texttt{region\_group\_South\_America} ($-1.42$) indicate that coffees from these regions are
much less likely to be classified as Specialty compared to the Caribbean/US reference.
\texttt{processing\_group\_Other} ($-1.17$) and \texttt{variety\_group\_Typica} ($-0.80$)
also showed strong negative effects. Defect counts reduce the likelihood of Specialty
classification, consistent with industry expectations.

% -----------------------------------------------
\section{Conclusion \& Business Implications}
% -----------------------------------------------

This project shows that production-related information---origin, variety, processing method,
altitude, and defect counts---can help predict coffee quality levels. The balanced logistic
regression model (Model~A) performs substantially better than a simple baseline, proving that it
is possible to estimate coffee quality from production-side features alone.

However, the gap between the production-feature model and the single sensory-score reference is
large. Model~A achieves a Macro-F1 of approximately 0.46 compared to approximately 0.83 for the
flavor-score model (Model~B), demonstrating that cupping data provides much stronger predictive
value than any combination of production features in this dataset.

\textbf{Business Implications:}
\begin{itemize}
  \item \textbf{For buyers and importers:} Coffee from East Africa, especially when semi-washed
        or naturally processed and with low defects, is often a strong indicator of high quality,
        allowing buyers to focus on the most promising samples and reduce cupping time and costs
        \citep{IburuCoffee2021, EbruCoffee2025}.
  \item \textbf{For farmers and cooperatives:} Lowering defect levels has the greatest effect on
        improving coffee grade. Better sorting and post-harvest processing tools can greatly increase
        the chance of reaching higher quality categories \citep{SCA2005, WEF2020}.
  \item \textbf{For quality graders and certifiers:} A simple screening model using basic origin
        and production data can reduce the need for full cupping while maintaining reliable
        classification results \citep{PDG2024}.
\end{itemize}

\textbf{Limitations:}
\begin{itemize}
  \item The dataset is small and has very few Commercial coffee samples, which may reduce model
        accuracy.
  \item Logistic regression assumes simple relationships; more advanced models might perform better.
  \item Some values (e.g., altitude) were estimated using country-level medians, potentially
        reducing accuracy.
\end{itemize}

\textbf{Future Work:}
\begin{itemize}
  \item Adding a small amount of sensory data could improve predictions considerably.
  \item Creating a real-time probability scoring system for coffee quality to support business
        decisions.
\end{itemize}

% -----------------------------------------------
% References
% -----------------------------------------------
\bibliographystyle{plainnat}
\begin{thebibliography}{99}

\bibitem[LeDoux(2018)]{LeDoux2018}
LeDoux, J. (2018). \textit{coffee-quality-database} [Data set]. GitHub.
\url{https://github.com/jldbc/coffee-quality-database}

\bibitem[SCA(2005)]{SCA2005}
Specialty Coffee Association. (2005). \textit{SCAA cupping protocols}. Specialty Coffee
Association of America.
\url{https://atlanticspecialtycoffee.com/wp-content/uploads/SCAA-Cupping-Protocols-2005.pdf}

\bibitem[Wintgens(2009)]{Wintgens2009}
Wintgens, J.~N. (Ed.). (2009). \textit{Coffee: Growing, processing, sustainable production---A
guidebook for growers, processors, traders and researchers} (2nd ed.). Wiley-VCH.

\bibitem[Iburu Coffee(2021)]{IburuCoffee2021}
Iburu Coffee. (2021, November 4). Coffee growing areas in Kenya: A comprehensive overview.
\textit{Iburu Coffee Blog}.
\url{https://iburucoffee.com/blogs/blog/coffee-growing-areas-in-kenya-a-comprehensive-overview}

\bibitem[Ebru Coffee Co.(2025)]{EbruCoffee2025}
Ebru Coffee Co. (2025). \textit{About us}. \url{https://ebrucoffeeco.com/pages/about-us}

\bibitem[Perfect Daily Grind(2024)]{PDG2024}
Perfect Daily Grind. (2024, July 23). Specialty coffee needs to find new ways to assess quality.
\textit{Perfect Daily Grind}.
\url{https://perfectdailygrind.com/2024/07/specialty-coffee-new-ways-to-assess-quality/}

\bibitem[WEF(2020)]{WEF2020}
World Economic Forum. (2020, August 3). How sales of certain types of coffee could transform
economies.
\url{https://www.weforum.org/stories/2020/08/quality-coffee-can-boost-local-economies-and-benefit-farmers/}

\bibitem[Harris et al.(2020)]{Harris2020}
Harris, C.~R., Millman, K.~J., van der Walt, S.~J., et al. (2020). Array programming with
NumPy. \textit{Nature}, \textit{585}(7825), 357--362.
\url{https://doi.org/10.1038/s41586-020-2649-2}

\bibitem[Hunter(2007)]{Hunter2007}
Hunter, J.~D. (2007). Matplotlib: A 2D graphics environment. \textit{Computing in Science \&
Engineering}, \textit{9}(3), 90--95. \url{https://doi.org/10.1109/MCSE.2007.55}

\bibitem[Pedregosa et al.(2011)]{Pedregosa2011}
Pedregosa, F., Varoquaux, G., Gramfort, A., et al. (2011). Scikit-learn: Machine learning in
Python. \textit{Journal of Machine Learning Research}, \textit{12}, 2825--2830.
\url{https://jmlr.org/papers/v12/pedregosa11a.html}

\bibitem[The pandas development team(2020)]{Pandas2020}
The pandas development team. (2020). \textit{pandas-dev/pandas: Pandas} (v1.0.5) [Software].
Zenodo. \url{https://doi.org/10.5281/zenodo.3509134}

\bibitem[Virtanen et al.(2020)]{Virtanen2020}
Virtanen, P., Gommers, R., Oliphant, T.~E., et al. (2020). SciPy 1.0: Fundamental algorithms
for scientific computing in Python. \textit{Nature Methods}, \textit{17}(3), 261--272.
\url{https://doi.org/10.1038/s41592-019-0686-2}

\bibitem[Waskom(2021)]{Waskom2021}
Waskom, M.~L. (2021). Seaborn: Statistical data visualization. \textit{Journal of Open Source
Software}, \textit{6}(60), 3021. \url{https://doi.org/10.21105/joss.03021}

\end{thebibliography}

% -----------------------------------------------
\section*{Appendix}
% -----------------------------------------------
\addcontentsline{toc}{section}{Appendix}

\begin{itemize}
  \item \textbf{Link to GitHub Repository:} \url{[Insert link]}
  \item \textbf{Data Source:} \url{https://github.com/jldbc/coffee-quality-database}
\end{itemize}

\end{document}
